\documentclass[11pt]{article}
\usepackage{fontspec}
\setmainfont{Times New Roman}
\setsansfont{Arial}
\setmonofont{Consolas}
\usepackage{amsmath,amssymb,mathtools}
\usepackage{geometry}
\usepackage{microtype}
\geometry{a4paper,margin=2.28cm}
\setlength{\parindent}{1.15em}
\setlength{\parskip}{0pt}
\providecommand{\mo}{\mathfrak{o}}
\providecommand{\mc}{\mathfrak{c}}
\emergencystretch=140em
\allowdisplaybreaks
\sloppy
\hbadness=10000
\vbadness=10000
\hfuzz=2pt
\vfuzz=10pt
\raggedbottom
\begin{document}
% Unit: Paper34 Section23 source-fidelity continuation v001.
% Language lane: Interslavic.
% Direct source: sources/paper34/source_fidelity/Noether_Paper34_Section23_ORIGINAL_SCAN_WITNESS_v001.tex.
\section*{Emmy Noether: Hiperkompleksne veličiny i teorija predstavjenj}
\emph{Papir 34, §23. Izvorno-točny prěvod iz skanovanogo německogo svědka.}

\subsection*{§ 23. Determinanta hiperkompleksnogo sistema}

Nehaj
\[
        \mo=a_1P+\cdots+a_hP
\]
bude hiperkompleksny sistem. Pridajemo k $P$ $h$ neodrěđene veličiny $x_1,\ldots,x_h$ i tvorimo
\[
        \mo^*=a_1P(x)+\cdots+a_hP(x).
\]
Neodrěđene veličiny $x_i$ imajut komutovati s $a_i$; tym už sut opreděljena pravila računa v $\mo^*$.

V $\mo^*$ leži „obči element“ sistema $\mo$:
\[
        w=a_1x_1+\cdots+a_hx_h.
\]
Ako v nekom predstavjenju $a_i\mapsto A_i$, togda elementu $w$ sootvětuje
\[
        W=A_1x_1+\cdots+A_hx_h.
\]
$W$ imenuje se sistemna matrica danogo predstavjenja; specialno, ako $a_i$ tvore grupu i $\mo$ zato jest grupovo kolco, to jest grupova matrica. Pri regularnom predstavjenju imajemo regularnu sistemnu matricu.

Elementy $w_{ij}$ matrice $W$ sut linearne formy v $x_i$. Zato sistemna determinanta $|W|$ jest stepena $n$, ako ide o predstavjenja $n$-togo stepena. Osoblivo regularna sistemna determinanta jest stepena $h$.

Sistemna determinanta ne měnja se pri prěhodu k ekvivalentnym predstavjenjam, bo
\[
        |PWP^{-1}|=|P|\,|W|\,|P^{-1}|=|W|.
\]
Pri prěhodu od bazy $(a_1,\ldots,a_h)$ k novoj bazě $(b_1,\ldots,b_h)$ i od $w=\sum a_ix_i$ k $w=\sum b_jy_j$, nove elementy $W$, i zato takože nova determinanta, dostavajut se iz staryh elementov substitucijeju
\[
        x_i=\sum_j\rho_{ij}y_j
\]
s regularnoju substitucijskoju matrico.

Ako dana jest kompozicijna serija modula predstavjenja, togda pri podobnom izboru bazy matrica $W$ ima formu
\[
        W=\begin{pmatrix}
        W_1&0&\cdots&0\\
        *&W_2&\cdots&0\\
        \vdots&\vdots&\ddots&\vdots\\
        *&*&\cdots&W_s
        \end{pmatrix}.
\]
Medžu determinantami $|W_i|$ libovoljnogo predstavjenja ne pojavjajut se nikake druge, než v regularnom predstavjenju sistema $\mo$, ili čak sistema $\mo/\mc$, kde $\mc$ jest radikal.

Ako $P$ jest algebraično zatvorjeno, togda determinanta $|W_i|$, k odnomu ireducibilnomu predstavjenju, jest ireducibilna funkcija v $x_i$, a k neekvivalentnym predstavjenjam prinadležet različne ireducibilne faktory.

Dokaz. Poneže vse ireducibilne predstavjenja $\mo$ sut takože predstavjenja $\mo/\mc$, možemo ograničiti se na kolco bez radikala $\mo/\mc$. V njem beremo za bazu matrične jedinice $c_{ik}^{(\nu)}$; togda obči element jest
\[
        w=\sum_{\nu,i,k}c_{ik}^{(\nu)}x_{ik}^{(\nu)}.
\]
Matrice ireducibilnyh predstavjenj sut
\[
        W_\nu=(x_{ik}^{(\nu)}).
\]
Funkcije $|W_\nu|=|x_{ik}^{(\nu)}|$ sut, kako poznano, ireducibilne i javno različne.

Da by izračunati $|W_i|$, možno razložiti regularnu sistemnu matricu $\mo/\mc$ na jej ireducibilne faktory. Vsaky ireducibilny faktor pojavja se toliko razov, koliky jest stepenj sootvětnogo ireducibilnogo predstavjenja, bo sootvětny ideal toliko razov pojavja se v kompozicijnoj seriji.

Možno takože izhoditi iz regularnoj sistemnoj matrice $\mo$, no togda vsaky ireducibilny faktor dostaje se čestěje, imenno toliko razov, koliko sootvětny prosty ideal pojavja se kako kompozicijny faktor medžu lěvymi idealami. V tom regularnom predstavjenju $\mo$ bylo smatrjano kako lěvy ideal; ako smatrjati $\mo$ kako pravy ideal, pojavja se druga regularna sistemna matrica, „antistrofna matrica“ pri Frobeniusu. Ona sodrži te same ireducibilne faktory, imenno sistemne determinanty vseh ireducibilnyh predstavjenj, ale možno s drugymi eksponentami; gl. primjer v § 10.

Sistemna determinanta komutativnogo sistema razpadaje se na linearne faktory, bo vse ireducibilne predstavjenja sut prvogo stepena. Te linearne faktory same sut ireducibilne predstavjenja, i zato pri abelovyh grupah davajut karaktere. Ta fakt byl izhodny punkt Dedekinda pri izučenju grupovoj determinanty neabelovyh grup.

\end{document}
